\documentclass[10pt,a4paper]{amsart}
\usepackage[margin=19mm]{geometry}
\usepackage{amsmath,amssymb,mathtools}
\usepackage[scr=rsfs,cal=euler]{mathalfa}
\usepackage{xfrac}
\usepackage[unicode,hidelinks,bookmarks=false]{hyperref}
\newcommand{\ie}{i.e.\ }
\newcommand{\cf}{cf.\ }
\newcommand{\resp}{respectively\ }
\newcommand{\Subsection}{\subsection}
\providecommand{\nobreakdash}{\nobreak\hbox{-}}
\setlength{\emergencystretch}{2em}
\multlinegap0pt
\usepackage{bm}
\usepackage{bbm}
\usepackage{dsfont}
\usepackage{xspace}
\usepackage{cancel}
\usepackage{mathtools}
\usepackage{amsthm}
\makeatletter
\@ifundefined{dcases}{\newtheorem{dcases}{Dcases}}{}
\@ifundefined{exer}{\newtheorem{exer}{Exercise}}{}
\@ifundefined{pr}{\newtheorem{pr}{Proposition}}{}
\@ifundefined{statement}{\newtheorem{statement}{Statement}}{}
\makeatother
\newcommand {\ii} {\mathrm{i}}
\title[Lectures on Gauge theories and Many-Body systems]{Lectures on Gauge theories and Many-Body systems}
\author{Igor Chaban, Nikita Nekrasov}
\date{Source: arXiv:2512.23099v2 (CC BY 4.0)}
\begin{document}
\maketitle

\noindent\emph{Source and licence.} Lectures on Gauge theories and Many-Body systems. Version arXiv:2512.23099v2 (CC BY 4.0), distributed under the Creative Commons Attribution 4.0 International licence, as recorded in the arXiv licence metadata for this exact version and shown on the abstract page https://arxiv.org/abs/2512.23099v2. Full rights notice: licenses/arxiv-2512.23099-CC-BY-4.0.txt.\par\smallskip

\noindent\textit{Editorial scope.} Counted unit: the Pr whose source label is \texttt{None} in the article (source0.tex lines 601--603, source label \texttt{None}), delivered together with the article's own complete original proof of it (source0.tex lines 604--628, one \texttt{proof} environment of 25 lines). No mathematics is written, repaired or re-derived: statement and proof are the article's own text line for line. Required context retained uncounted: none. Every definition, display and label of those lines is retained unchanged. Omitted from this package and not redistributed: 1--600, 629--2777. Nothing inside a retained excerpt is omitted or altered: every retained body is delivered exactly as the article's lines are written, so no omission receipt of any kind accompanies this package. Citations: the retained excerpts carry no citation command at all, so no bibliography entry of the article is transcribed and no reference is redistributed. Reference adaptation: none was needed and none was made; every reference inside the delivered unit resolves inside it, so no unresolved reference or citation is produced. Printed slips kept literal: no printed word, symbol, hypothesis or number of the counted statement or of its proof was altered. Layout and class: the article's own class is \texttt{article} with packages including inputenc, babel, graphicx, geometry, mathtools, amsmath, cancel, slashed, misccorr, siunitx, amsthm, amssymb, wrapfig, booktabs; the wrapper is the lane's pinned \texttt{amsart} editorial wrapper at 10pt on A4, which restates the theorem-like environments the retained text needs and adds the article's own macro definitions that the retained text needs (\texttt{\textbackslash ii}), with extra wrapper packages (bm, bbm, dsfont, xspace, cancel, mathtools). The delivered numbering is the unit's own. Assets: no retained span contains a figure environment, a graphics-inclusion command, a PDF-page inclusion, a table or a TikZ picture; no image file is copied into this package, the package declares no asset dependency and no image file is redistributed with it. Licence: version arXiv:2512.23099v2 of this article is distributed under the Creative Commons Attribution 4.0 International licence, as recorded in the arXiv licence metadata for this exact version and shown on the abstract page https://arxiv.org/abs/2512.23099v2. A full rights notice is delivered at licenses/arxiv-2512.23099-CC-BY-4.0.txt. Characters: every retained excerpt is pure ASCII, so the delivered engine, XeLaTeX with its default Latin Modern text font, reports no missing character.
\begin{pr}
    $\mathcal{P}_{S}$ is symplectomorphic to $T^{*}\{(x_1 , ..., x_n ): x_1 < ... < x_n \}$.
\end{pr}

\begin{proof}
Let us smoothly pick a representative in every $U(N)-$orbit of $\mu^{-1}(0)^{ss}$. The set of representatives is a leaf transversal to $U(N)-$orbits and diffeomorphic to the quotient space $\mathcal{P}_{S}$. The functions $p_1, ..., p_N$ are thus defined on $\mathcal{P}_{S}$ in terms of restriction on the leaf. The choice of representatives (gauging) is our freedom. To simplify the calculations one should use this freedom wisely. We propose the following choice. In every $U(N)-$orbit we pick a diagonal $X = \mathrm{diag}(x_1, ..., x_N ), ~ x_1 < ... < x_N$.
The remaining symmetry is $\mathrm{Stab}_{U(N)}(X) = U(1) \times ... \times U(1) \subset U(N)$.
In the $U(1)^{\times N}-$orbit we choose $z_i \in \mathbb{R}_{+}$. This fixes the gauge completely. In fact due to the constraint $\mu = 0$ the coordinates $P_{ij}$ and $z_i$ are determined completely with a given $x_i$ and $p_i$ on the leaf chosen.
\begin{statement}{Problem 2}\normalfont\label{prob2}
Show that the conditions $\mu = 0$, $X = \mathrm{diag}(x_1, ..., x_n)$, $x_1 < ... <x_N$ , $z_i \in \mathbb{R}_{+}$ imply
\begin{align*}
\begin{dcases}
        z_i = \sqrt{\nu}, \;i = 1, ..., N\\\\
        P_{ij} = \frac{{\ii}\nu}{x_i - x_j}, \;i \neq j,
    \end{dcases} 
\end{align*}
the matrix $P$ hence has the form
\begin{align}\label{Pmatrix}
    P_{ij} = p_{i} \delta_{ij} + \frac{{\ii}\nu (1 - \delta_{ij})}{x_i - x_j}.
\end{align}
\end{statement}
Thus, identifying $p_i$ with a coordinates on the cotangent fibers we see that the leaf is diffeomorphic to $T^{*}\{(x_1 , ..., x_n ): x_1 < ... < x_n \}$. The following exercise identifies the symplectic forms and thereby concludes the proposition.
\begin{exer}\normalfont
    Using the defining property $\pi^{*}\omega_{S} = i^{*}\omega_{L}$ for $\omega_{S}$, where $i: \mu^{-1}(0)^{ss} \hookrightarrow \mathcal{P}_{L}$ and $\pi: \mu^{-1}(0)^{ss} \rightarrow \mu^{-1}(0)^{ss}/U(N)$, check that
    \begin{align*}
        \omega_{S} = \sum\limits_{i = 1}^{N} d p_{i} \wedge d x_{i}
    \end{align*}
\end{exer}
\end{proof}

\end{document}
